После указанных исправлений все три дополнительные проверки завершены
и приняты. Для уменьшенного зазора получены пороги 15,523 и
19,563 $\uJ$ для боковых продольного и радиального каналов.
При уменьшении $c$ до 11,875 нм они равны 13,883 и 19,558,
а при увеличении $a=b$ до 4,725 нм --- 13,583 и
18,854 $\uJ$. Во всех трёх случаях продольный канал
остаётся эффективнее радиального. Коэффициенты материала номинальной
формы при окончательном повторе малого зазора точно совпадают
с исходными; изменён только пространственный порядок.

\begin{table}[!htbp]\centering\small
\caption{Чувствительность порогов двух ведущих боковых каналов.
Каждый параметр изменяется отдельно, остальные сохраняются.
Флюенсы приведены в $\uJ$; при вариациях формы $g=0{,}5$ нм.
$E_q$ --- энергия перехода КТ, несущая остаётся 2,042 эВ.}
\label{tab:sensitivity}
\begin{tabular}{lrr}\toprule
Вариация & Боковой продольный & Боковой радиальный\\\midrule
Номинальные параметры & 16,075 & 20,229\\
$g=0{,}3$ нм & 15,523 & 19,563\\
$g=0{,}7$ нм & 16,659 & 20,961\\
$c=11{,}875$ нм & 13,883 & 19,558\\
$c=13{,}125$ нм & 18,156 & 21,102\\
$a=b=4{,}275$ нм & 18,740 & 22,092\\
$a=b=4{,}725$ нм & 13,583 & 18,854\\
$E_q-5$ мэВ & 15,757 & 20,532\\
$E_q+5$ мэВ & 16,756 & 20,382\\
$\gamma_1-20\%$ & 16,069 & 20,222\\
$\gamma_1+20\%$ & 16,081 & 20,236\\
$\gamma_\varphi-20\%$ & 16,124 & 20,707\\
$\gamma_\varphi+20\%$ & 16,074 & 19,953\\
\bottomrule\end{tabular}

\end{table}

Объединённая проверка всех 12 однофакторных вариаций принимает
сохранение порядка двух ведущих каналов
(табл.~\ref{tab:sensitivity}). Изменения формы оказывают наибольшее
влияние: продольный боковой порог меняется от 13,583 до
18,740 $\uJ$, то есть от $-15{,}50$ до $+16{,}58$\%
относительно номинального. Для изменения $g$ на $\pm0{,}2$ нм
изменения составляют $-3{,}44$ и $+3{,}63$\%.
Изменение энергии перехода на $\pm5$ мэВ даёт пороги
15,757 и 16,756 $\uJ$. Вариации $\gamma_1$ на
$\pm20$\% влияют значительно слабее, менее чем на 0,04\%
для ведущего канала, что согласуется с большой величиной $T_1$
относительно времени считывания. Вариация чистого дефазирования
на те же 20\% меняет продольный порог менее чем на 0,31\%,
радиальный --- до 2,36\%. Такое различие отражает зависимость
послеимпульсной эволюции от геометрии обратной связи.

\figone{validation_sensitivity.png}{.98}{Однофакторная чувствительность
после восстановления трёх первоначальных отказов. Показано максимальное
по контрольной группе абсолютное относительное изменение порога;
группа включает два ведущих канала и изолированную КТ.
Знак изменения можно установить по табл.~\ref{tab:sensitivity}.
Изменение несущей рассматривается отдельно.}{fig:sensitivity}

При оценке разделимости используется консервативный относительный
запас $\eta=5$\%, не являющийся статистическим доверительным
интервалом. На исходной сетке точки продольного бокового канала
$g=0{,}5;\ 0{,}75;\ 1$ нм имеют перекрывающиеся с лучшей точкой
интервалы $\Fth(1\pm\eta)$ и образуют область близкой эффективности.
Поэтому численно обоснован диапазон 0,5--1 нм на данной сетке,
а не единственный точно локализованный оптимальный зазор.
Контрольная точка $g=0{,}3$ нм относится к проверке чувствительности
и не включается задним числом в исходную сетку выбора.

Принятый критерий имеет ограниченную область действия: исходные
дискретные кандидаты плюс поочерёдные вариации параметров двух
выбранных каналов. Он не исключает появления другого лидера
среди всех пяти каналов при совместном изменении параметров,
не представляет статистический анализ ансамбля частиц и не
снимает физических ограничений субнанометровых зазоров.
